\begin{table*}[t]
\centering
\small
\caption{Parte II: resultados dos algoritmos no G13.}
\label{tab:algoritmos}
\begin{adjustbox}{max width=\linewidth}
\begin{tabular}{lp{0.78\linewidth}}
\toprule
Algoritmo & Resultado \\
\midrule
BFS & origem = askreddit; alcança 34.671 vértices; excentricidade = 6 \\
DFS & origem = askreddit; visita 34.671 vértices \\
Euler (não dirigido) & 25.306 vértices de grau ímpar; conexo: não; euleriano: não; semieuleriano: não \\
Euler (dirigido) & 32.257 vértices com entrada ≠ saída; fracamente conexo: não; euleriano: não \\
Dijkstra & w = 1/count; maior distância a partir de askreddit = 4,59 \\
Bellman-Ford & mesmas distâncias que Dijkstra: sim; rodadas até convergir: 6 (10 origens: 8–10); ciclo negativo: não \\
Floyd-Warshall & só em subgrafos (n ≤ 2000): O(n³) e matriz 34.671² ≈ 1,2·10⁹ células (\textasciitilde{}9,6 GB) inviáveis \\
Tarjan (SCC) & 24.071 componentes fortemente conexas; maior = 11.564 vértices (32,3\%); 23.949 unitárias; 497 fracamente conexas \\
Tarjan (pontes e articulações) & maior componente: 18.057 pontes (14,6\% das arestas) e 5.006 vértices de articulação (14,4\%) \\
Kruskal & 34.670 arestas; peso total = 25.435,531 \\
Prim & 34.670 arestas; peso total = 25.435,531 (igual ao Kruskal: sim) \\
\bottomrule
\end{tabular}
\end{adjustbox}
\end{table*}
